\subsubsection*{Introduction}
Generally, prosthetics users rely on visual feedback and secondary cues, such as sound or changes in pressure on the residual limb \cite{schofield_applications_2014}.
Sensory feedback, particularly vibrotactile, is associated with better object recognition, manipulation and grip force control \cite{schiefer_artificial_2018}\cite{rosenbaum-chou_development_2016}.
While sensory feedback is often delivered through electrical stimulation for more traditional prostheses, we hypothesise that Intraosseous Transcutaneous Amputation Prostheses(ITAP) can be the vehicle for transmission of haptic feedback.
Haptics can take the form of thermal feedback, force and much more --- in particular we focus on vibrotactile, i.e. the touch perception of vibrations.
However, little is understood about the consequence of vibrations on the prosthesis-skin integration or on its attachment to bone.

\begin{figure}[h]\centering
    \incfig[]{ITAP}
    \caption{ITAP nomenclature for its parts, context of its insertion to the body and the areas of interest for measurement of displacement.}
    \label{fig:ITAP}
\end{figure}
